% Part: incompleteness
% Chapter: introduction
% Section: historical-background

\documentclass[../../../include/open-logic-section]{subfiles}

\begin{document}

\olfileid{inc}{int}{bgr}

\olsection{చారిత్రక నేపథ్యం}

ఈ విభాగంలో అసంపూర్ణతా సిద్ధాంతాలను వాటి సందర్భంలో ఉంచడానికి తోడ్పడే చారిత్రక
పరిణామాలను క్లుప్తంగా చర్చిస్తాం. ముఖ్యంగా, గణిత తర్క చరిత్రను అత్యంత సంక్షిప్తంగా
సమీక్షించి, తరువాత గణిత పునాదుల చరిత్ర గురించి కొన్ని మాటలు చెబుతాం.

\begin{digress}
  ``గణిత తర్కం'' అనే పదబంధానికి రెండు అర్థాలున్నాయి. ``గణితంలోని తర్కం''
  అన్నప్పుడు మాదిరిగా, ``గణిత'' అనే పదాన్ని విషయాన్ని వర్ణించేదిగా భావిస్తే అది
  గణిత హేతుచింతన సూత్రాలను సూచిస్తుంది; ``తర్కపు గణితం'' అన్నప్పుడు మాదిరిగా,
  పద్ధతులను వర్ణించేదిగా భావిస్తే అది హేతుచింతన సూత్రాల గణిత అధ్యయనాన్ని
  సూచిస్తుంది. కింది వివరణలో గణిత తర్కం తరచూ ఒకేసారి ఈ రెండు అర్థాల్లోనూ ఉంటుంది.
\end{digress}

తర్క అధ్యయనం మౌలికంగా క్రీ.పూ. సుమారు 384--322లో జీవించిన అరిస్టాటిల్‌తో
ప్రారంభమైంది. ఆయన రచనలు \emph{వర్గాలు}, \emph{పూర్వ విశ్లేషణలు},
\emph{ఉత్తర విశ్లేషణలు} శాస్త్రీయ హేతుచింతన సూత్రాలను వ్యవస్థాబద్ధంగా
అధ్యయనం చేస్తాయి; వాటిలో త్రివాక్య నిగమనాన్ని సమగ్రంగా, క్రమబద్ధంగా పరిశీలించడం
కూడా ఉంది.

మధ్యయుగాలంతటా పాండిత్య తత్వశాస్త్రంపై అరిస్టాటిల్ తర్కమే ఆధిపత్యం చెలాయించింది;
పద్దెనిమిదవ శతాబ్దంలో కూడా కాంట్, అరిస్టాటిల్ తర్కం పరిపూర్ణమనీ దానికి సవరణ
అవసరం లేదనీ వాదించాడు. కానీ త్రివాక్య నిగమన సిద్ధాంతం గణిత హేతుచింతనలోని అత్యంత
ఉపరితల అంశాలకంటే మరేదీ నమూనీకరించలేనంత పరిమితమైనది. అంతకు ఒక శతాబ్దం ముందు,
న్యూటన్ సమకాలికుడైన లైబ్నిట్స్ తార్కిక హేతుచింతనకు సంపూర్ణ ``కలనం''ను
ఊహించాడు; అటువంటి కలనాన్ని రూపొందించే దిశగా కొన్ని ప్రాథమిక అడుగులు వేస్తూ,
మౌలికంగా ప్రతిజ్ఞావాక్యాత్మక తర్కపు ఒక రూపాన్ని వర్ణించాడు.

పందొమ్మిదవ శతాబ్దం తర్క చరిత్రలో కీలక మలుపు. 1854లో జార్జ్ బూల్
\emph{ఆలోచనా నియమాలు} రచించాడు; అందులో ఆధునిక వివరణలకు చాలా దగ్గరగా ఉండే
ప్రతిజ్ఞావాక్యాత్మక తర్కపు సమగ్ర బీజగణిత అధ్యయనం ఉంది. 1879లో గాట్‌లోబ్ ఫ్రేగె
తన \emph{బెగ్రిఫ్‌ష్రిఫ్ట్} (భావలిపి)ను ప్రచురించాడు. అది ప్రతిజ్ఞావాక్యాత్మక
తర్కాన్ని పరిమాణకాలతో, సంబంధాలతో విస్తరించి మొదటిస్థాయి తర్కాన్ని కలిగి ఉంటుంది.
నిజానికి ఫ్రేగె తార్కిక వ్యవస్థల్లో ఉన్నత-స్థాయి తర్కమూ, మరికొన్ని అంశాలూ ఉన్నాయి.
తన \emph{అంకగణితపు మౌలిక నియమాలు}లో, అంకగణితమంతటినీ పూర్తిగా తార్కికమైన
పరికల్పనల నుండి తన బెగ్రిఫ్‌ష్రిఫ్ట్‌లో వ్యుత్పాదించవచ్చని చూపడానికి ఫ్రేగె
ప్రయత్నించాడు. దురదృష్టవశాత్తూ, 1902లో రసెల్ చూపినట్లు ఈ పరికల్పనలు వైరుధ్యమైనవి.
అయితే ఆ వైరుధ్య స్వీకృతాన్ని పక్కనపెడితే, ఫ్రేగె దాదాపు ఒంటిచేత్తో ఆధునిక తర్కాన్ని
ఆవిష్కరించాడు---ఇది ఆశ్చర్యకరమైన సాధన. బూల్ తరువాత బీజగణిత దృక్పథం గల పియర్స్,
ష్రోడర్ తదితరులు కూడా పరిమాణక తర్కాన్ని స్వతంత్రంగా అభివృద్ధి చేశారు.
\sourcecorrection{OLTEINCINT-001}{ఆంగ్ల మూలంలోని ``purely logical assumption'' అనే ఏకవచనాన్ని, వెంటనే బహువచనంగా సూచించిన పరికల్పనలకూ ఫ్రేగె వ్యవస్థలోని పలు తార్కిక పరికల్పనలకూ అనుగుణంగా ``పూర్తిగా తార్కికమైన పరికల్పనలు''గా ఇచ్చాం.}

ఇప్పుడు గణిత పునాదులలోని పరిణామాలవైపు తిరుగుదాం. తర్కం గణితంలో ముఖ్య పాత్ర
పోషిస్తుంది కనుక, ఇప్పుడే వివరించిన పరిణామాలతో దీనికి గణనీయమైన పరస్పర సంబంధం
ఉండటం సహజమే. ఉదాహరణకు, గణితమంతటినీ తన తార్కిక చట్రంపైనే ఆధారపరచవచ్చని
చూపాలనే స్పష్టమైన ఉద్దేశంతో ఫ్రేగె తన తర్కాన్ని అభివృద్ధి చేశాడు. ముఖ్యంగా,
కాంట్ వాదించినట్లుగా గణితం పూర్వానుభవ \emph{సంశ్లేషణాత్మక} సత్యాలతో కాక,
పూర్వానుభవ \emph{విశ్లేషణాత్మక} సత్యాలతో కూడినదని చూపాలని అతను ఆశించాడు.

గణితం సిసలుగా గ్రీకులతో పుట్టిందని చాలామంది భావిస్తారు. క్రీ.పూ. సుమారు 300లో
రచించిన యూక్లిడ్ \emph{ఎలిమెంట్స్}, నిష్కర్షకూ ఖచ్చితత్వానికీ ఇచ్చిన ప్రాధాన్యంతో,
అప్పటికే పరిపక్వ గ్రీకు గణితానికి ప్రతినిధి. యూక్లిడ్ \emph{ఎలిమెంట్స్}లోని
నిర్వచనాలు, నిరూపణలు నేటి ఉన్నత పాఠశాల జ్యామితి పాఠ్యపుస్తకాల్లో దాదాపు యథాతథంగా
కొనసాగుతున్నాయి (ఉన్నత పాఠశాలల్లో జ్యామితి ఇంకా ఎంతవరకు బోధిస్తారో అంతవరకు).
ఈ గణిత హేతుచింతన నమూనా గణితంలోనే కాక తత్వశాస్త్ర శాఖల్లోనూ నిష్కర్షక వాదనకు
ఆదర్శంగా పరిగణించబడింది. (స్పినోజా నైతిక, మతపర వాదనలను కూడా యూక్లిడ్ శైలిలో
ప్రస్తావించాడు; చూడటానికి అది వింతగా ఉంటుంది!)

పదిహేడవ శతాబ్దంలో న్యూటన్, లైబ్నిట్స్ కలనాన్ని ఆవిష్కరించారు. (దీనిపై శతాబ్దాల
పాటు తీవ్రమైన ప్రాథమ్య వివాదం సాగింది; కానీ నేడు చాలామంది పండితులు ఈ రెండు
అభివృద్ధులు ఎక్కువగా స్వతంత్రమైనవని భావిస్తారు.) కలనంలో, ఉదాహరణకు, అనంతంగా
చిన్న పరిమాణాల అనంత మొత్తాలపై హేతుచింతన ఉంటుంది. ఈ లక్షణాలు బిషప్ బర్క్లీ
విమర్శలకు దారితీశాయి; దేవునిపై నమ్మకం తన కాలపు గణితంకంటే తక్కువ హేతుబద్ధమైనది
కాదని అతను వాదించాడు. పద్దెనిమిదవ శతాబ్దంలో కలన పద్ధతులు విస్తృతంగా వాడబడ్డాయి;
ఉదాహరణకు, అనంత మొత్తాలతో లెక్కలను ఉపయోగించి లియోనార్డ్ ఆయిలర్ అద్భుత ఫలితాలు
సాధించాడు.

పందొమ్మిదవ శతాబ్దంలో గణితజ్ఞులు కలనానికి మరింత దృఢమైన పునాది కల్పించి బర్క్లీ
విమర్శలకు సమాధానం ఇవ్వడానికి ప్రయత్నించారు. కోషీ, వైయర్‌స్ట్రాస్, బోల్జానో
తదితరుల కృషి వల్ల, అతిసూక్ష్మ పరిమాణాల ప్రస్తావనే లేకుండా ``ఎప్సిలాన్‌లు,
డెల్టాలు'' పరంగా హద్దులు, అవిచ్ఛిన్నత, అవకలనం, సమాకలనం గురించి మన సమకాలీన
నిర్వచనలు ఏర్పడ్డాయి. అదే శతాబ్దంలో తరువాత, వాస్తవ సంఖ్యలతో సహా కలనంలోని అన్ని
అంశాలను సహజ సంఖ్యల పరంగా వివరించడానికి గణితజ్ఞులు ఇంకా ముందుకు వెళ్లారు.
(క్రోనెకర్: ``పూర్ణ సంఖ్యలను దేవుడు సృష్టించాడు; మిగతాదంతా మానవుని కృషి.'')
1872లో డెడెకిండ్ ``అవిచ్ఛిన్నత మరియు అకరణీయ సంఖ్యలు'' రాసి, అకరణీయ సంఖ్యలను
సహజ సంఖ్యల జతలుగా చూడగల పరిమేయ సంఖ్యల సమితులుగా తీసుకుని వాస్తవ సంఖ్యలను ఎలా
``నిర్మించవచ్చో'' చూపాడు. 1888లో అతను ``వాస్ జిండ్ ఉండ్ వాస్ జొల్లెన్ డీ
త్సాలెన్'' (సుమారుగా, ``సహజ సంఖ్యలు ఏమిటి, అవి ఎలా ఉండాలి?'') రాసి, సహజ
సంఖ్యలను పూర్తిగా ``తార్కిక'' పరంగా వివరించడానికి ప్రయత్నించాడు. 1887లో
క్రోనెకర్ ``యూబర్ డెన్ త్సాల్‌బెగ్రిఫ్'' (``సంఖ్య భావన గురించి'') రాసి, అన్ని
గణిత వస్తువులనూ పూర్ణ సంఖ్యల పరంగా సూచించడం గురించి చెప్పాడు; 1889లో జుజెప్పె
పియానో సహజ సంఖ్యలకు సాంకేతిక, సంకేతాత్మక స్వీకృతాలను ఇచ్చాడు.
\sourcecorrection{OLTEINCINT-002}{ఆంగ్ల మూలంలోని ``all mathematical object''లో బహువచన విభక్తి లోపించింది. క్రోనెకర్ ఉద్దేశించిన సాధారణ వాదనకు అనుగుణంగా ``అన్ని గణిత వస్తువులు'' అనే అర్థాన్ని ఇచ్చాం.}

పందొమ్మిదవ శతాబ్దం చివర్లో అనంతాన్ని వ్యవహరించడంలో కొత్త ధైర్యం కూడా వచ్చింది.
అంతకుముందు అనంత వస్తువులనూ నిర్మాణాలనూ (సహజ సంఖ్యల సమితి వంటివాటిని) చాలా
జాగ్రత్తగా పరిగణించేవారు; ``అనంతంగా అనేకం'' అంటే ``ఎన్ని కావాలంటే అన్ని'' అని,
``హద్దులో సమీపిస్తుంది'' అంటే ``ఎంత దగ్గర కావాలంటే అంత దగ్గర అవుతుంది'' అని
అర్థం చేసుకునేవారు. కానీ అనంతాన్ని ఉన్నదున్నట్లుగానే స్వీకరించవచ్చని జార్జ్
కాంటర్ చూపాడు. కాంటర్, డెడెకిండ్ తదితరుల కృషి, నేడు విస్తృతంగా ఆమోదించబడిన
గణితపు సాధారణ సమితి-సిద్ధాంత అవగాహనను ప్రవేశపెట్టడానికి తోడ్పడింది.
\sourcecorrection{OLTEINCINT-003}{ఆంగ్ల మూలంలోని ``Work ... help to introduce''లో కాలసమ్మతి లోపించింది. పూర్తయిన చారిత్రక పరిణామాన్ని సూచించేలా ``కృషి ... తోడ్పడింది'' అని ఇచ్చాం.}

ఇలా మనం ఇరవయ్యవ శతాబ్దపు తర్క, పునాది పరిణామాలకు చేరుకుంటాం. 1902లో ఫ్రేగె
తార్కిక వ్యవస్థలోని విరోధాభాసాన్ని రసెల్ కనుగొన్నాడు. 1904లో జెర్మెలో
``ఎంపిక స్వీకృతం''గా పిలిచే దాన్ని ఉపయోగించి కాంటర్ సుక్రమీకరణ సూత్రాన్ని
నిరూపించాడు; ఈ స్వీకృతపు సమర్థనీయతపై చాలా చర్చ జరిగింది. 1910 నుండి 1913 మధ్య
రసెల్, వైట్‌హెడ్‌ల \emph{ప్రిన్సిపియా మాథమాటికా} మూడు సంపుటాలు వెలువడి,
గణితాన్ని తార్కిక పునాదులపై స్థాపించాలనే ఫ్రేగె కార్యక్రమాన్ని విస్తరించాయి.
దురదృష్టవశాత్తూ, పూర్తిగా తార్కికమైనవిగా సమర్థించడం కష్టంగా కనిపించిన రెండు
సూత్రాలను---అనంతత్వ స్వీకృతాన్నీ ``కుంచనీయత'' స్వీకృతాన్నీ---రసెల్, వైట్‌హెడ్
అంగీకరించాల్సి వచ్చింది. 1900లలో గణితంలోని ``అపూర్వాపేక్ష నిర్వచనాల'' వాడకాన్ని
పొయంకారే విమర్శించాడు; 1910లలో, బహిష్కృత మధ్యమ సూత్రం ($!A \lor \lnot !A$)
వాడకాన్ని నివారించే ``అంతఃప్రజ్ఞావాద'' పునాదిపై గణితమంతటినీ తిరిగి స్థాపించాలని
బ్రౌవర్ ప్రతిపాదించడం మొదలుపెట్టాడు.

నిజంగా అవి విచిత్రమైన రోజులు!{} గణితమంతటినీ తర్కానికి కుదించే కార్యక్రమాన్ని
ఇప్పుడు ``తర్కవాదం'' అంటారు; పైన చెప్పిన కష్టాల వల్ల అది విఫలమైందనే అభిప్రాయం
సాధారణంగా ఉంది. అంతఃప్రజ్ఞావాద మానసిక నిర్మాణాల పరంగా గణితాన్ని అభివృద్ధి చేసే
కార్యక్రమాన్ని ``అంతఃప్రజ్ఞావాదం'' అంటారు; అది దైనందిన గణితంపై అతిగా కఠినమైన
పరిమితులు విధిస్తుందని భావిస్తారు. శతాబ్దం మారే సమయానికి, సర్వకాల అత్యంత ప్రభావశీల
గణితజ్ఞుల్లో ఒకడైన డేవిడ్ హిల్బర్ట్, కాంటర్, డెడెకిండ్ ప్రవేశపెట్టిన కొత్త నైరూప్య
పద్ధతులకు బలమైన మద్దతుదారు: ``కాంటర్ మనకోసం సృష్టించిన స్వర్గం నుండి మనల్ని
ఎవరూ తరిమివేయలేరు.'' అదే సమయంలో, ఈ కొత్త పద్ధతులపై (వింతగా, ఇప్పుడు
``సాంప్రదాయిక'' అని పిలిచే పద్ధతులపై) వచ్చిన పునాది విమర్శలను అతను గంభీరంగా
తీసుకున్నాడు. రెండు ప్రయోజనాలనూ పొందే మార్గాన్ని అతను ప్రతిపాదించాడు:
\begin{enumerate}
\item సాంప్రదాయిక పద్ధతులను సాంకేతిక స్వీకృతాలు, నియమాలతో సూచించాలి; గణిత
  ప్రశ్నలను ఒక స్వీకృత వ్యవస్థలోని \tetoken{సూత్రాలుగా}{!!{formula}s} సూచించాలి.
\item ఈ సాంకేతిక నిగమన వ్యవస్థలు అవైరుధ్యమైనవని నిరూపించడానికి సురక్షితమైన,
  ``పరిమితీయ'' పద్ధతులను ఉపయోగించాలి.
\end{enumerate}

మొదటి లక్ష్యాన్ని సాధించే దిశగా హిల్బర్ట్ కృషి చాలా ముందుకు వెళ్లింది. 1899లో
తన ప్రసిద్ధ గ్రంథం \emph{జ్యామితి పునాదులు}లో జ్యామితికి అతను ఇదే పని చేశాడు.
తరువాతి సంవత్సరాల్లో, హిల్బర్ట్ జ్యామితికి చేసిన పనిని గణితంలోని ఇతర రంగాలకు
చేయడానికి అతనూ అతని విద్యార్థులు, సహచరుల్లో పలువురూ కృషి చేశారు. అంకగణితానికీ
విశ్లేషణకూ హిల్బర్టే స్వీకృత వ్యవస్థలను ఇచ్చాడు. జెర్మెలో సమితి సిద్ధాంతానికి
స్వీకృతీకరణ ఇచ్చాడు; దాన్ని ఫ్రెంకెల్, స్కోలెమ్, ఫాన్ నాయ్‌మన్ తదితరులు
విస్తరించారు. 1920ల మధ్య నాటికి, గణిత ``మొత్తానికి'' స్వీకృతీకరణ అనే పేరును
కోరిన రెండు విధానాలు ఉన్నాయి: రసెల్, వైట్‌హెడ్‌ల \emph{ప్రిన్సిపియా
మాథమాటికా}, జెర్మెలో--ఫ్రెంకెల్ సమితి సిద్ధాంతంగా ప్రసిద్ధమైన వ్యవస్థ.

1921లో ఈ వ్యవస్థలు అవైరుధ్యమైనవని నిరూపించే లక్ష్యంతో హిల్బర్ట్ ఒక పరిశోధనా
కార్యక్రమాన్ని ప్రారంభించాడు. ఈ కార్యక్రమంలో అతని విద్యార్థుల్లో పలువురు,
ప్రత్యేకించి బెర్నాయ్స్, అకెర్మాన్, తరువాత గెంట్సెన్ అతనికి సహకరించారు. ఈ లక్ష్యాన్ని
సాధించే మౌలిక ఆలోచన, గణితంలో వైరుధ్యపు \tetoken{వ్యుత్పత్తి}{!!{derivation}}
సాధ్యమా అనే ప్రశ్నను సంకేతాల సాధ్య క్రమాల గురించిన సంయోజనాత్మక సమస్యగా మార్చడం.
అంటే అంకగణితం, విశ్లేషణ లేదా సమితి సిద్ధాంతానికి స్వీకృత వ్యవస్థలోని స్వీకృతాల
నుండి, ఉదాహరణకు, $!A \land \lnot !A$కు సరైన \tetoken{వ్యుత్పత్తి}{!!{derivation}}
అయ్యే వాక్య క్రమాలు సాధ్యమా అని ప్రశ్నించడం. అటువంటి సంకేత క్రమం అసాధ్యమని
నిరూపించే నిరూపణ కూడా గణిత నిరూపణే కనుక, ఈ స్వీకృత వ్యవస్థల్లో సాంకేతికంగా
వ్యక్తీకరించదగినదై ఉంటుంది. మరో మాటలో, ఉదాహరణకు అంకగణితం అవైరుధ్యమైనదని చెప్పే
ఏదో వాక్యం~$\OCon$ ఉంటుంది. పైగా దాని నిరూపణకు చాలా పరిమితమైన ``పరిమితీయ''
పద్ధతులు మాత్రమే అవసరమైతే, ఆ వాక్యం ఆయా వ్యవస్థల్లో నిరూపించదగినదై ఉండాలి.

అభివృద్ధి చేసిన స్వీకృత వ్యవస్థలు ప్రతి గణిత ప్రశ్ననూ పరిష్కరించాలనే రెండవ
లక్ష్యాన్ని రెండు విధాలుగా ఖచ్చితంగా చెప్పవచ్చు. ఒక విధంగా ఇలా చెప్పవచ్చు: గణితపు
స్వీకృత వ్యవస్థ భాషలోని ప్రతి వాక్యం~$!A$కు, $!A$ గానీ $\lnot !A$ గానీ
స్వీకృతాల నుండి నిరూపించదగినదై ఉండాలి. ఇది నిజమైతే, స్వీకృతాల ఆధారంగా నిరూపించనూ
ఖండించనూ వీలులేని వాక్యాలు ఉండవు; స్వీకృతాలు పరిష్కరించని ప్రశ్నలు ఉండవు. ఈ
ధర్మం గల స్వీకృత వ్యవస్థను \emph{సంపూర్ణం} అంటాం. ఏదైనా నిర్దిష్ట వాక్యానికి
ఈ రెండు ప్రత్యామ్నాయాల్లో ఏది నిజమో కనుగొనడం కష్టమైన పని కావచ్చు. కానీ సూత్రతః
దాన్ని చేసే పద్ధతి ఉండాలి. నిజానికి, హిల్బర్ట్ పరిగణించిన స్వీకృత,
\tetoken{వ్యుత్పత్తి}{!!{derivation}}
వ్యవస్థలకు సంపూర్ణత ఉంటే, అటువంటి పద్ధతి ఉంటుందని అనుసరిస్తుంది---అయితే హిల్బర్ట్
దీన్ని గుర్తించలేదు. ప్రశ్నను అర్థం చేసుకునే రెండవ విధానం మరింత బలమైన షరతు:
ఇచ్చిన వాక్యం~$!A$ స్వీకృతాల నుండి వ్యుత్పాదించదగినదా కాదా అని నిర్ణయించే
యాంత్రిక, గణనాత్మక పద్ధతి ఉండాలి.

1931లో గ్యోడెల్ రెండు ``అసంపూర్ణతా సిద్ధాంతాలను'' నిరూపించి, తగినంత బలమైన,
అవైరుధ్యమైన, ప్రభావకంగా స్వీకృతీకరించిన గణిత వ్యవస్థ సంపూర్ణం కాలేదనీ, అలాంటి
వ్యవస్థ తన స్వంత అవైరుధ్యాన్ని వ్యక్తం చేసే వాక్యాన్ని సాధారణంగా నిరూపించలేదనీ
చూపాడు.\sourcecorrection{OLTEINCINT-004}{ఆంగ్ల మూలం ఏ స్వీకృత వ్యవస్థా సంపూర్ణం కాదని పరిమితుల్లేకుండా చెప్పి, అవైరుధ్య వాక్యం నిరూపించనూ ఖండించనూ వీలుకాదని రెండవ అసంపూర్ణతా సిద్ధాంతానికి మించిన వాదన చేసింది. తరువాతి అధ్యాయాల ఖచ్చిత ప్రకటనలకు అనుగుణంగా, ప్రభావక స్వీకృతీకరణ, అవైరుధ్యం, తగిన బలం అనే షరతులను చేర్చి, నిరూపించలేమనే సరైన ఫలితానికే వాక్యాన్ని పరిమితం చేశాం.}

ఇది హిల్బర్ట్ అసలు కార్యక్రమానికి ప్రాణాంతక దెబ్బ ఇచ్చింది. అయితే గణితంలో తరచూ
జరిగినట్లే, ఇది ఉత్తేజకరమైన కొత్త పరిశోధనా మార్గాలను కూడా తెరిచింది. గణితమంతటినీ
ఆవరించే ఒకే సాంకేతిక వ్యవస్థ లేకపోతే, మరింత పరిమిత పరిధిగల వ్యవస్థలను అభివృద్ధి
చేసి వాటిలో ఏమి నిరూపించవచ్చో పరిశీలించడం సమంజసం. ఈ వ్యవస్థల అవైరుధ్యాన్ని
స్థాపించడానికి తక్కువ పరిమితులున్న నిరూపణ పద్ధతులను అభివృద్ధి చేయడం, వాటి అవైరుధ్యాన్ని
నిరూపించడం ఎంత కష్టమో కొలిచే మార్గాలను కనుగొనడం కూడా సమంజసం. (దాదాపు) ప్రతి
సాంకేతిక వ్యవస్థ పరిష్కరించలేని ప్రశ్నలు కలిగి ఉంటుందని గ్యోడెల్ చూపాడు కనుక,
ఒక నిర్దిష్ట సాంకేతిక వ్యవస్థ పరిష్కరించలేని ``ఆసక్తికర'' ప్రశ్నలను వెతకడం,
వాటిని పరిష్కరించడానికి ఆ వ్యవస్థ ఎంత బలంగా ఉండాలో కనుగొనడం సమంజసం. నేటికీ
తర్కవేత్తలు ఈ ప్రశ్నలను నిరూపణ సిద్ధాంతం అనే కొత్త గణిత శాఖలో పరిశోధిస్తున్నారు.

\end{document}
